\cotefile{93}
\scanpage{179}
\unsure{vient|vent|revient} équivalemment d'exprimer le th.~3, la relation
\[
\boxed{\operatorname{res}\binom{\widetilde C\omega}{\widetilde f_1,\ldots,\widetilde f_n}
 =\operatorname{res}\binom{\omega}{f_1^p\ldots f_n^p}}
\]
Lorsque la dim relative est 1, cette formule s'écrit
\[
\begin{aligned}
\operatorname{res}\frac{\widetilde C\omega}{\widetilde f}
 &=\operatorname{res}\frac{\omega}{f^p}
 &\text{i.e.}\quad &\boxed{\operatorname{res}\widetilde C\unsure{\varpi|\overline\omega}=\operatorname{res}\varpi}\\
\overset{\rotatebox{90}{$=$}}{\inserted{\operatorname{res}\widetilde C\left(\frac{\omega}{f}\right)}}
 &&&\text{si on pose }\varpi=\omega/\unsure{\widetilde f|f|f^p}.
\end{aligned}
\]
et lorsque de plus $S$ est parfait i.e. $F_S:S\longrightarrow S$ un isom, de sorte que $C$ est défini, on peut l'écrire
\[
\begin{aligned}
\deleted{\left(\operatorname{res}\frac{C\omega}{f}\right)^p=\operatorname{res}\frac{\omega}{f^p}}
 &\qquad \boxed{\left(\operatorname{res} C\varpi\right)^p=\operatorname{res}\varpi}\\
\deleted{\overset{\rotatebox{90}{$=$}}{\operatorname{res} C\left(\frac{\omega}{f^p}\right)}}&
\end{aligned}
\]

\uline{Remarques} \textcircled{1} Il y a aussi un hom de Cartier $f_i:f_*(\Omega^{\deleted{\unsure{k}}i}_{X/S})\longrightarrow\Omega^i_{X^{p/S}/S}$, défini grâce à l'hom \inserted{de Cartier} pour $i=n$ et la formule de projection. Mais comme $f^\ast:\Omega^i_{X^{p/S}/S}\longrightarrow f_*\Omega^i_{X/S}$ est nul pour $i\ne0$, on conclut que $f_i=0$ si $i\ne n$. En particulier si $i\ne0$, l'hom $f_i$ ne coïncide \uline{pas} en général avec l'opération de Cartier, \inserted{sur les formes fermées}.

\textcircled{2} Il y a lieu de définir aussi l'opération de Cartier \deleted{\unsure{disons}} sur les modules dualisants relatif $\underline{\omega}_{X/S}$, \inserted{sans hyp. de lissité,} lorsque la formation de celui-ci commute au chgt. de base $S\xrightarrow{F_S}S$ \unsure{.} (p.ex. $S$ parfait i.e. $F_S$ un isom), \deleted{\unsure{et}} $X$ étant loc de type fini sur $S$ loc \unsure{noeth} disons. On trouve
\[
 C=f_{X/S\,*}:f_{X/S\,*}(\omega_{X/S})\longrightarrow
 \omega_{X^{p/S}/S}\simeq\widetilde\omega_{X/S}
\]
d'où, en appliquant $f_{X/S}^{-1\,*}$, \inserted{et composant avec $f_{X/S}'(\widetilde\omega_{X/S})\longrightarrow\widetilde\omega_{X/S}^{(p)}$,} un hom (pas linéaire, mais $\Lambda$-linéaire)
\[
 \omega_{X/S}\longrightarrow\omega_{X/S}^{(p)}
\]
% archive sheet number: 178
\scanpage{180}
Appliquons le th.~3 à la dualité globale, en supposant $X$ \uline{propre et lisse} sur $S$. Soit $E$ module loc. libre sur $X$ (ou complexe parfait $\ldots$), $\widetilde E$ son image inverse sur $X^{p/S}$, \inserted{dont l'image inverse sur $X$ par $f_{X/S}$ est $E^{(p)}$.} Alors on a les homomorphismes
\[
\deleted{\left[H^i(X,E)\xrightarrow{p\text{-lin}}\right]H^i(X^{p/S},\widetilde E)\longrightarrow H^i(X,E^{(p)})}
\]
dit ``puissance p.ième''
\[
 \underset{\substack{\rotatebox{90}{$\simeq$}\\\mathbb R f_*(E)^{(p)}}}{\mathbb R f_*(\widetilde E)}
 \xrightarrow{\unsure{r|\sigma|\varphi}}\mathbb R f_*(E^{(p)})\tag*{($\ast$)}
\]
qui n'est autre essentiellement que l'hom de changemt de base associé au carré commutatif

\diagram{\begin{tikzcd}[ampersand replacement=\&]
{}^{E}X \arrow[r,no head,"F_X"] \arrow[d,no head,"f"] \& X^{F_X^\ast(E)=E^{(p)}}\arrow[d,no head]\\
S\arrow[r,no head,"F_S"]\&S
\end{tikzcd}}

Considérons le \uline{transposé} de $(\ast)$, nous savons que c'est l'hom
\[
\begin{aligned}
&\hspace{10em}\simeq Rg_*\bigl(Rf_{X/S\,*}(E^{(p)}\otimes\Omega^n_{X/S})\bigr)\\
&Rf_*(\check E^{(p)}\otimes\Omega^n_{X/S}[n])\mathrel{\text{---}}
 Rf_*(\check{\widetilde E}\otimes\Omega^n_{\widetilde X/S}[n])
\end{aligned}
\]
associé à l'hom trace
\[
 f_{X/S\,*}(\check E^{(p)}\otimes\Omega^n_{X/S})
 =\check{\widetilde E}\otimes f_{X/S\,*}(\Omega^n_{X/S})
 \longrightarrow\check{\widetilde{\unsure{E|\mathbb E}}}\otimes\Omega^n_{X^{p/S}/S}\tag*{($\ast\ast$)}
\]
déduit, par tensorisation avec $\check{\widetilde E}[n]$ \deleted{\unsure{avec}} de l'hom trace habituel. Celui-ci à son tour s'interprète comme l'hom \deleted{trace} de Cartier.

En particulier, on voit que le transposé de
\[
 \underset{\substack{\rotatebox{90}{$\simeq$}\\R^nf_*(E)^{(p)}}}{R^nf_*(\widetilde E)}
 \longrightarrow R^nf_*(E^{(p)})
\]
\deleted{[supposons $R^nf_*(E)$ et $R^nf_*(E^{(p)})$ loc. libres, pour simplifier]}
% archive sheet number: 179
\scanpage{181}
est l'hom
\[
\begin{aligned}
f_*(\check E^{(p)}\otimes\Omega^n_{X/S})&\longrightarrow
 g_*(\check{\widetilde E}\otimes\Omega^n_{X^{p/S}/S})\\
\deleted{f_*(\check E\otimes\Omega}&
 \quad\rotatebox{90}{$\simeq$}\uparrow\ \text{\unsure{ss|sous} réserve de commutation à $F_S$ de $f_*(\check E\otimes\Omega^n_{X/S})$}\\
&g_*(\check E\otimes\Omega^n_{X/S})^{(p)}
\end{aligned}
\]
donné par l'opération de Cartier $(\ast\ast)$ ci-dessus.

Le cas le plus important est celui où $E=\underline O_X$, on trouve que l'hom de puissance p.ième
\[
 R^nf_*(\underline O_X)^{(p)}\mathrel{\text{---}}R^nf_*(\underline O_X)
\]
est transposée de l'opération de Cartier
\[
 \widetilde C:f_*(\Omega^n_{X/S})\mathrel{\text{---}}g_*(\Omega^n_{X^{p/S}/S})
 \mathrel{\underset{\text{ss. réserve}}{\overset{\sim}{\longleftarrow}}}
 f_*(\Omega^n_{X/S})^{\otimes p}
\]
$\overline{\mathrm M}$ énoncé pour les $R^if_*(\underline O_{X\deleted{\unsure{{}^{p/S}}}})^{(p)}\mathrel{\text{---}}R^if_*(\underline O_X)$, \inserted{sous réserve que c'est $R^if_*(\underline O_{X^{p/S}/S})$} à cela près qu'il faut prendre l'hom sur $R^{n-\unsure{1|i}}f_*(\Omega^n_{X/S})\longrightarrow R^{n-i}f_*(\Omega^n_{X/S})^{\unsure{\otimes p|\otimes1/p|(p)}}$ \inserted{[sous réserve]} déduit de $\widetilde C$ par application de $R^{n-i}f_*$.

Un hom intéressant est
\[
 R^1f_*(\underline O_{X^{p/S}})\longrightarrow R^1f_*(\underline O_X)
\]
dans le cas où $f$ est propre, \deleted{car} \deleted{\unsure{c'est}} il s'identifie à l'hom. des sections pour les fibrés vectoriels
\[
 \underset{\substack{\parallel\\(\underline{\operatorname{Lie}}\,\underline{\operatorname{Pic}}_{X/S})^{(p/S)}}}
 {\underline{\operatorname{Lie}}\,\underline{\operatorname{Pic}}_{X^{p/S}/S}}
 \longrightarrow\underline{\operatorname{Lie}}\,\underline{\operatorname{Pic}}_{X/S}
\]

\margin{\uline{À vérifier} (et y donner un sens sous des conditions aussi générales que possible)}
% slanted
Nous admettrons ici que \deleted{\unsure{l'}} l'hom \deleted{\unsure{sur Pic}}

\diagram{\begin{tikzcd}[ampersand replacement=\&]
\underline{\operatorname{Pic}}_{X^{p/S}/S}\arrow[r,no head]\arrow[d,no head,"\rotatebox{90}{$\simeq$}" description,draw=none]\&\underline{\operatorname{Pic}}_{X/S}\\
\underline{\operatorname{Pic}}_{X/S}^{p/S}\arrow[ur,"V"]\&
\end{tikzcd}}
\margin{Il est immédiat en tout cas que $F_{\mathrm{Pic}}V=p\,\mathrm{id}$, $VF_{\mathrm{Pic}}=p.\mathrm{id}$ \deleted{\unsure{ce qui}}}
est le \uline{Verschiebung} , d'où il résulte que l'hom.
% archive sheet number: 180
\scanpage{182}
en question induit sur \uline{Lie} est l'opération de puissance p.ème (est-ce dans le livre de Demazure-Gabriel ?). \margin{?} L'hom transposé est donc donné par l'opération de Cartier sur les formes différentielles.

\textcircled{7} \uline{Cas \inserted{lisse} de \deleted{\unsure{la}} dim relative un : un accouplement remarquable.}

\deleted{Soit} Plus généralement, au lieu de dim relative 1,
\[
 \deleted{Q=f_{X/S\,*}(\underline O_X)/\underline O_{X^{p/S}}}
\]
supposons que $X$ sur $X^{p/S}$ \deleted{\unsure{soit de loc une p-base,}} \inserted{soit de \unsure{t.f.|r.f.} et} $\Omega^1_{X/S}$ $(=\Omega^1_{X/X^{p/S}})$ soit \uline{libre de rang $1$}, ce qui signifie aussi que $X$ a loc. une $p$-base à 1 élément sur $S$. Posons
\[
\begin{aligned}
 Q&=f_{X/S\,*}(\underline O_X)/\operatorname{Im}\underline O_{X^{p/S}}
 =f_{X/S\,*}(\underline O_X)/\Lambda\\
 &\simeq\underline O_X/f^{-1}\underline O_S[\unsure{\underline o|\underline O}_X^p]
\end{aligned}
\]
qui est un module loc libre de rang $p-1$ sur $S$. On va définir un \uline{accouplement \inserted{$\Lambda$-bilinéaire} canonique}
\[
 \begin{aligned}
 Q\times Q&\longrightarrow\deleted{\unsure{\Lambda}}\,\Omega^1_{X^{p/S}/S}
 \otimes_{\underline O_{X^{p/S}}}\Lambda\\
 &\qquad(=\Omega^1_{X^{p/S}/S}\quad\text{si p.ex. $X/S$ lisse})
 \end{aligned}\tag*{($\ast$)}
\]
\margin{$-$} \deleted{ou} \deleted{\unsure{pour|plutôt}} ce qui revient au m\textsuperscript{e}, un hom $\Lambda$-bil.
\[
 \varphi:f_{X/S\,*}(\underline O_X)\times f_{X/S\,*}(\underline O_X)
 \longrightarrow\Omega^1_{X^{p/S}/S}\otimes_{\underline O_{X^{p/S}}}\Lambda
\]
qui est nul sur $\Lambda\times f_{X/S\,*}(\underline O_X)$ et sur $f_{X/S\,*}(\underline O_X)\times\Lambda$.
% archive sheet number: 181
\scanpage{183}
On prend
\[
 \varphi(f,g)=\widetilde C(f\,dg)
\]

a) $\Lambda$-bilinéaire : clair

b) nul si \deleted{$f$ section de $\Lambda$} $f=1$ ou $g=1$ : clair, car $\widetilde C$ est nul sur les diff. exactes

c) $\varphi$ est antisymétrique, i.e. $\left\{\begin{array}{l}\text{alterné si }p\ne2\\\text{symétrique si }p=2\end{array}\right.$

i.e. $\varphi(f,g)+\varphi(g,f)=0$ i.e. $\widetilde C(f\,dg+g\,df)=0$, en effet $f\,dg+g\,df=d(fg)=0$.

\uline{Proposition} L'accouplement $(\ast)$ est une dualité parfaite de $\Lambda$-modules.

On peut supposer qu'il existe une $p$-base relative $\{x\}$, et on a une base de $Q$ sur $1$ formé par les classes des \deleted{$x,x^{\unsure{2|3}},\ldots,x^{p-1}$}. $x^i$ $1\le i\le p-1$. On a
\[
 \widetilde C(x^i\,dx^j)=\widetilde C(jx^{i+j-1}\,dx)=
 \begin{cases}0&\text{si }i+j\ne p-1\\j&\text{si }i+j=p-1\end{cases}
\]
donc la matrice de la forme $\varphi$ est
\[
 \begin{pmatrix}
 &0&&p-1\\
 &&\reflectbox{$\ddots$}&\\
 &2&&0\\
 1&&&
 \end{pmatrix},
\]
son déterminant, au signe près, est donc $1.2.\ldots(p-1)\in\mathbb F_p^\ast$, \margin{$\Delta=-1\in\mathbb F_p^\ast$} donc inversible, O.K.

\uline{Corollaire} \deleted{$\det Q$} On a un isom canonique
\[
 Q\xrightarrow{\sim}\check Q\otimes(\Omega^1_{X^{p/S}/S}\otimes\Lambda)
\]
$\Downarrow$

d'où un isom canonique
\[
 \det Q\simeq\det\check Q\otimes(\Omega^1_{X^{p/S}/S}{}^{\otimes(p-1)}\otimes\Lambda)
\]
% archive sheet number: 182
\scanpage{184}
et par suite un isom canonique
\[
 \det(Q)^{\otimes2}\simeq{\Omega^1_{X^{p/S}/S}}^{\otimes(p-1)}\otimes\Lambda
 \tag*{($\ast\ast$)}
\]

\uline{Cas} 1 \quad $p=2$, on trouve
\[
 \deleted{\det Q^{\otimes}\simeq}\ \Omega^1_{X^{p/S}/S}\otimes\Lambda\simeq\det(Q)^{\otimes2}
\]
donc $\Omega^1_{X^{p/S}/S}\otimes\Lambda$ est \uline{canoniquement un carré}. Si $S$ est parfait, \inserted{et $X$ réduit, i.e. $\Lambda\simeq\underline O_{X^{p/S}}$,} ceci implique que $\Omega^1_{X/S}$ est lui-m\textsuperscript{e} canoniquement un carré. \deleted{Je doute que \unsure{ce|ceci} pour $X$} \deleted{\unsure{affine}} \deleted{(cf. résultat de)} Ceci n'est pas vrai si $S=\operatorname{Spec}k$, $k$ \inserted{corps} imparfait de car. 2, \inserted{m\textsuperscript{e} si $X$ lisse sur $k$,} comme \deleted{\unsure{et}} on voit \deleted{\unsure{déjà}} par l'exemple de Serre-Tate-\unsure{Fröhlich} : il suffit de prendre l'exemple d'une courbe \inserted{lisse propre} rationnelle sans point -- \margin{?} \deleted{\illegible{1}} il y en a sur tout corps imparfait de car 2

\uline{Cas} 2 \quad $p\ne2$, donc $\frac{p-1}{2}$ est un entier. On aimerait alors préciser $(\ast\ast)$ en un isom can
\[
 \psi:\det(Q)\simeq{\Omega^1_{X^{p/S}/S}}^{\otimes(\frac{p-1}{2})}\otimes\Lambda
 \left[\simeq{\Omega^1_{X/S}}^{\otimes(\frac{p-1}{2})}\otimes_{\underline O_X}\Lambda\right]
\]
[Sinon, il y aurait un \unsure{revêtement} de groupe $\mu_2\simeq\mathbb Z/2\mathbb Z$ de $X^{p/S}$, canoniquement associé à la situation]. Reprenons la base de $Q$, $(\dot x^i)_{1\le i\le p-1}$, \deleted{\unsure{et soit|et comme}} d'où une base de \deleted{\unsure{la}} $\det Q$, $\dot x\wedge\dot x^2\wedge-\wedge\dot x^{p-1}$, et la valeur de son carré, par l'isom $(\ast\ast)$, est $\Delta(d\widetilde x^{\otimes(p-1)}\otimes\Lambda)$, où $\Delta$ est le déterminant de la matrice écrite plus haut, savoir $(-1)^{(p-1)-1}1.2.\ldots(p-1)=\deleted{\unsure{+}}-(1.2.\ldots(p-1))$. Comme $\mathbb F_p^\ast$ est un groupe cyclique d'ordre \deleted{\unsure{un}} pair \unsure{pour|$p-1$} le produit de ses éléments est $-1$, donc on trouve $\Delta=\unsure{+1|1}$
% archive sheet number: 183
\scanpage{185}
Ceci suggère \deleted{d'associer à $\dot x$} de définir $\unsure{\varphi|\psi}$ par
\[
 \boxed{\psi(\dot x\wedge\dot x^2\wedge\text{--}\wedge\dot x^{p-1})
 =d\widetilde x^{\otimes\frac{p-1}{2}}\otimes 1_\Lambda}
 \tag*{($\ast\ast\ast$)}
\]
Ce qui donne une racine carrée de l'isom $(\ast\ast)$, et dont \unsure{revient|restent|restait} à prouver que $\psi$ est indépendant du choix de la $p$-base $\{x\}$. \deleted{Soit $y$ une autre $p$-base,}
\[
 \deleted{y=a_{p-1}x^{p-1}\mathrel{\text{---}}+a_0\qquad(a_i\in\Lambda)}
\]
\deleted{d'où}
\[
 \deleted{dy=[(p-1)a_{p-1}x^{p-2}\mathrel{\text{---}}+a_1]\,dx,}
\]
\deleted{donc on a}
\[
 \deleted{(p-1)a_{p-1}x^{\unsure{p-1|p-2}}\mathrel{\text{---}}+a_1\in{\unsure{\mathbb B|B}}^\ast}
\]
Or, le passage de l'un à l'autre se fait par un signe. On trouve un \deleted{\unsure{hom}} \inserted{morphisme} du schéma \inserted{$P$} [sur \deleted{\unsure{$\Lambda$}} \inserted{$Z=(X^{p/S},\Lambda)$}] des $p$-bases relatives, dans le revêtemt trivial à deux feuillets $(\mathbb Z/2\mathbb Z)_Z$ de $Z$. \deleted{\unsure{Cette ap}} Or \deleted{\unsure{$Z$}} $P$ sur $Z$ est à fibres connexes, et on voit donc que \deleted{cette application est} ce morphisme est \unsure{const\textsuperscript{ant}|const.|constante} fibre à fibre, ce qui donne ce qu'on veut $\ldots$.

Application : la dualité en cohomologie globale ($X$ lisse propre sur $S$ de dim rel. 1.)

\margin{?} Généralisation des relations précédentes à la dim supérieure ?
% archive sheet number: 184
\scanpage{186}
Cas $X\mathrel{\overset{f}{\text{---}}}S$ propre \inserted{lisse} de dim relative 1. On a un diagramme autodual [dualité = symétrie p.r. au centre)

\diagram{\begin{tikzcd}[ampersand replacement=\&,row sep=small,column sep=small]
\&\&0\&\&\\
\&\& R^1f_*(\underline O_{X/S})\arrow[u]\&\&\\
0\arrow[r]\& R^1f_*(\underline O_X)^{(p)}\arrow[r,"\alpha"]\arrow[ur,"F\ (\mathrm{frob})"]\&\mathcal H^1_{\mathrm{DR}}\arrow[r,"\beta"]\arrow[u]\&f_*(\Omega^1_{X/S})^{(p)}\arrow[r]\&0\\
\&\& f_*(\Omega^1_{X/S})\arrow[u,no head]\arrow[ur,"C\ (\text{op. de Cartier})"']\&\&\\
\&\&0\arrow[u]\&\&
\end{tikzcd}}
\margin{\uline{ligne et colonne \unsure{ex.ct|ex.cte}}}

\margin{$\ast$) $\mathcal H^1_{\mathrm{DR}}=\mathcal H^1_{\mathrm{DR}}(X/S)=R^1f_*(\Omega^\bullet_{X/S})$}

\margin{$\ast\ast$) Colonne est la filtration de $\mathcal H^1$ par $\mathcal H^{1,0}$ et $\mathcal H^{0,1}$}

\margin{$\ast\ast\ast$) $\alpha$ se déduit par passage au quotient de $F_{\mathcal H^1}:\mathcal H^{1(p)}\mathrel{\text{---}}\mathcal H^1$}

\margin{$\ast\ast\ast\ast$) $\beta={}^{t}\alpha$, déduit de $ {}^{t}F_{\mathcal H^1}$ par restriction}

\margin{$\ast\ast\ast\ast\ast$) $F$ et $C$ sont \deleted{duaux} \inserted{transposés} l'un de l'autre, ce sont la puiss. p.ième et Cartier.}

\deleted{On a $\beta\alpha=0$, et on a équivalence entre les conditions}

\deleted{(i) $\alpha$ injectif}

\deleted{(ii) $\beta$ surjectif}

\deleted{(iii) la suite $0\mathrel{\text{---}}R^1f_*(\underline O_X)^{(p)}\xrightarrow{\alpha}\mathcal H^1_{\mathrm{DR}}\xrightarrow{\beta}f_*(\Omega^1_{X/S})^{(p)}\longrightarrow0$ est exacte}

\deleted{(iii bis) $\operatorname{Ker}\beta=\operatorname{Im}\alpha$ \quad i.e. $\operatorname{Ker}\beta=(\operatorname{Ker}\beta)^\perp$ (on a tjs $\supset$) \quad i.e. $\operatorname{Im}\alpha=(\operatorname{Im}\alpha)^\perp$ (on a tjs $\subset$)}

\deleted{(iii ter) $\dim\operatorname{Ker}\beta=\dim\operatorname{Im}\alpha$}

\deleted{(iv) $\operatorname{rang}\alpha=g$ \quad i.e. $\operatorname{rang}\beta=g$ \quad i.e. $\operatorname{rg}F_{\mathcal H^1}=g$}

\margin{\deleted{\uline{NB} on a tjrs $\left\{\begin{aligned}\operatorname{Ker}\alpha&=(\operatorname{Im}\beta)^\perp\\\operatorname{Im}\alpha&=(\operatorname{Ker}\beta)^\perp\end{aligned}\right.$ \deleted{\illegible{2}}}}
\deleted{Pour ceci, il suffit que}
% archive sheet number: 185
\scanpage{187}
% Manuscript upside down in the scan; reading order after a half-turn.
\diagram{\begin{tikzcd}[ampersand replacement=\&,execute at end picture={\draw (current bounding box.south west)--(current bounding box.north east);}]
\deleted{X}\arrow[r,"{\deleted{f_{\unsure{\varnothing|0}}^{(i)}}}"]\&\deleted{X^{[i]}}\\
\deleted{T_0}\arrow[u,"{\deleted{\varphi_0}}"]\arrow[r,"{\deleted{i}}"']\&\deleted{T}\arrow[u,"{\deleted{u}}"']
\end{tikzcd}}
\deleted{tels que $\left\{\begin{array}{l}\varphi_0\in\unsure{\mathscr S|\mathcal S}_0\\i\in\mathcal N\end{array}\right.$}

\deleted{Considérons \deleted{\unsure{le cas}} plus gén. une flèche d'imm. $\delta:X\longrightarrow Z$, \inserted{de $(\mathrm{Sch})_{/S}$} et la catégorie analogue formée des}
\diagram{\begin{tikzcd}[ampersand replacement=\&,execute at end picture={\draw (current bounding box.south west)--(current bounding box.north east);}]
\deleted{X}\arrow[r,"{\deleted{\delta}}"]\&\deleted{Z}\\
\deleted{T_0}\arrow[u,"{\deleted{\varphi_{\unsure{0|u}}}}"]\arrow[r,"{\deleted{i}}"]\&\deleted{T}\arrow[u,"{\deleted{u}}"']
\end{tikzcd}}
\deleted{que par abus de langage nous allons désigner par $\unsure{C|\mathcal C}_{/\delta}$. Si $Z_{(n)}$ est le n.ième voisinage infinitésimal de $X$ dans $Z$, \deleted{alors} et $\delta_n:X\longrightarrow Z_{(n)}$, alors $C_{/\delta}$ est la \uline{Lim} des sous-catégories $C_{/\delta_n}$, qui cette fois \deleted{\unsure{aura}} [si \deleted{\unsure{$\delta_n$}} du moins \deleted{\unsure{$\delta_n$|$X$}} $Z_n\in\operatorname{Ob}\mathscr S$] \deleted{et que $X\hookrightarrow Z_n$ est \unsure{un hom.|une hom.} sur $X$ \unsure{(que|par} $(C)$ une rétraction sur $X$]},}
% archive sheet number: 186
\scanpage{188}
\uline{Connexions et stratifications} \uuline{intégrables} : (Suite, problèmes divers)
% archive sheet number: 187
\scanpage{189}
$X/S$ \quad schéma sur $S$.

1) Un module \inserted{qu. coh.} $E$ sur $X$, avec connexion rel/$S$, est dit \deleted{\unsure{le|de}} \uline{trivial} si $E\simeq f^\ast(F)$, $F$ un module \inserted{qu. coh.} sur $S$ [isom respectant connexion]. \margin{Il serait plus naturel à certains égards de travailler avec les PD-connexions [c'est pareil dans cas lisse]} \uline{NB} Si $X\longrightarrow S$ est \deleted{\unsure{couvrant fppf|connexe fpqc}} \inserted{lisse surjectif} et à fibres géométriquement connexes, \inserted{$S$ de carac 0, alors,} le foncteur $F\longrightarrow f^\ast(F)$, des modules qu. coh. sur $S$ dans les modules qu. coh. sur $X$ à connexion rel/$S$, est pleinement fidèle [du moins si $F$ de présentation finie].

On dit que $E$ est \uline{loc. trivial} ou \uline{Zar-intégrable} si $E$ est trivial localm\textsuperscript{t} sur $(X\longrightarrow S)$ pour la topologie de Zariski. On dit que $E$ est \uline{intégrable}, si la m\textsuperscript{e} condition est satisfaite, loc. pour top. étale.

2) Supposons $S$ de car $p>0$, et que $X$ sur $S$ admette loc. une $p$-base. Alors \deleted{$E$ est loc. trivial} conditions suivantes sont équivalentes sur $E$ qu coh.

a) $E$ est \uline{Zar. intégrable}

b) $E$ est intégrable

c) $E$ est à courbure et $p$-courbure nulle.

d) $E$ est \uline{trivial} relativement à $X\longrightarrow(X^{p/S},\underline O_{X^{p/S}}/J\simeq\Lambda\deleted{\unsure{\underline O_S|{}_{\!S}}})$

De plus, le foncteur $F\mathrel{\text{---}}{f'}_{X/S}^\ast(F)$ des Modules cohérents sur $X^{p/S\,\prime}=(X^{p/S},\underline O_{X^{p/S}}/J=\Lambda)$ aux modules \inserted{qu. coh. sur $X$} à connexion intégrable rel/$S$, est une équivalence de catégories.
% archive sheet number: 188
\scanpage{190}
3) Soit $G$ un groupe \deleted{fini} étale sur $S$, et considérons un $G_X$-torseur $P$ sur $X$ (à droite). \deleted{Pour toute} \inserted{Considérons une} représentation de $G$ \deleted{dans} un module qu. coh. $F$ sur $S$. Alors $G_X$ opère sur $F_X$, et on constate qu'il opère en laissant invariant la connexion canonique de $F_X$. Comme les Modules à connexion intégrable \inserted{rel/$S$} sur des objets du site étale de $X$ forment évidemment un champ, on voit que $G_X$ \inserted{(considéré comme \inserted{faisceau en} groupes sur $S$)} opère sur \deleted{\unsure{chaque|cha}} section \inserted{$F_X$,} du dit \deleted{\unsure{$F$|$E$}} champ, qu'on peut donc tordre par le $G_X$-torseur $P$. On trouve donc un module à connexion intégrable
\[
 P\mathbin{\times^{G_X}} F_X\; .
\]
\margin{En fait, on trouve m\textsuperscript{e}, par le m\textsuperscript{e} argument, une \uline{stratification} \uline{intégrable} sur $P\mathbin{\times^{G_X}}F_X$. Cela revient au m\textsuperscript{e} si $X/S$ lisse et $S$ de car. 0.}

Cette façon de construire des \deleted{connexions} \inserted{stratifications} intégrables est \deleted{\unsure{une construction|une connexion} \illegible{1}} universelle dans une large mesure. \margin{\unsure{Jouer|Prouver|Donnons} stabilité par sous-modules, quotients etc} \deleted{Supposons que $X\longrightarrow S$ soit plat, \deleted{\unsure{et loc.}} \inserted{\deleted{\unsure{loc.}} de prés. finie,} à fibres géométriques connexes, et qu'il ait une section $e$. Supposons} Du moins si $S$ est réduit au spectre d'un corps $k$, et que $X$ est p. ex. connexe et \deleted{de} loc. de présentation finie sur $k$, muni d'un pt $x\in X$ \deleted{\unsure{rationnel sur}} \inserted{\unsure{résiduel|rationnel} sur} $k$, on peut former le pro-groupe \inserted{strict} \deleted{\unsure{par}} (pas néc. profini si $X$ n'est pas normal) $\underline{\pi}_1(X,x)$, et la catégorie des Modules \inserted{\uline{coh.}} à stratification intégrables sur $X$ doit être équivalente à celle des représentations linéaires de dim finie de $\underline{\pi}_1(X,x)$. On doit pouvoir généraliser au cas d'une base quelconque $S$, en se
% archive sheet number: 189
\scanpage{191}
donnant un paquet $\mathbb L$ de nbs premiers, et supposons que l'on a un \deleted{pro}groupe \inserted{profini} $\underline{\pi}_1^{\mathbb L}(X,e)$, relativement à une section $e$ de $X/S$ [plus généralement, un sous-schéma $\unsure{\underline e|\underline E}$ de $X$ fini, \deleted{et} radiciel surjectif sur $S$] Je suppose $X\mathrel{\text{---}}S$ plat de présentation finie, \inserted{à fibres géom\textsuperscript{t} connexes,} $S$ loc. noeth $(?)$ : la \inserted{cat. des} représentations linéaires de $\underline{\pi}_1^{\mathbb L}(X,e)$ dans des modules cohérents sur $S$ doit être équivalente à celle des \deleted{représentations} Modules coh. sur $X$ à strat. intégrables rel/$S$, qui sur chaque fibre \inserted{géométrique} correspondent à une représentation du groupe fondamental qui \deleted{soit} se fasse via un $\mathbb L$-groupe.

Plus généralement, si on a un \deleted{groupe \unsure{profini}} pro-(groupe \deleted{\unsure{étale|fini}} \inserted{étale}) \deleted{\unsure{loc. const.|loc. constant}} strict $\Pi$ sur $S$, et un pro-torseur $P$ sous $\Pi_X$ sur $X$, tel que \inserted{sur chaque} fibre géom \deleted{\unsure{$f_{\bar s}$|fibre}} $\Pi_{\bar s}$ apparaisse comme \uline{quotient} du groupe fondamental élargi de $X_{\bar s}$, alors on \deleted{doit} trouve une \margin{?} équivalence entre la catégorie des repr. (continues) de $\Pi$ \unsure{ds|dans} des modules coh. sur $S$, et celle des Modules coh. à stratifications intégrables sur $X$ qui, sur chaque fibre géom, veulent bien être subordonnées à $\Pi_{\bar s}$ (regardé comme quotient du groupe fondamental de $X_{\bar s}$,. \ldots
% archive sheet number: 190
\scanpage{192}
4) Soit $X$ de type fini sur un corps, $X$ \deleted{\unsure{géom.}} \uuline{géom.} \uuline{unibranche}, on veut étudier les modules inversibles sur $X$ à stratification intégrable, \deleted{\unsure{indécomposable|de descente|nécessaire}}. \margin{loc. de t.f. suffit ?}

Soit $\underline L$ un module \inserted{inversible} stratifié sur $X$. Je dis que la stratification est intégrable ssi $\exists$ entier $n>0$, \inserted{premier à la car.,} tel que $\underline L^{\otimes n}$ soit trivial comme module stratifié [Le ``si'' est valable sans hyp. géom. unibranche], i.e ssi \deleted{on trouve} \inserted{$\underline L$ donne} un élément d'ordre fini \inserted{premier à $p$} dans le groupe des modules inversibles stratifiés. A fortiori, une \inserted{première} condition nécessaire est qu'il existe un entier $n>0$ \inserted{premier à $p$} tel que $\underline L^{\otimes n}$ admette une section partout $\ne 0$, i.e. que $\underline L$ donne élément d'ordre fini \inserted{premier à $p$} dans le $\operatorname{Pic}(X)$.

Cette condition étant remplie, il reste à exprimer que $\underline L^{\otimes n}$ est à \deleted{connexion} \inserted{stratification} \uline{intégrable}. On est donc ramené à caractériser \deleted{\unsure{les}} parmi les \deleted{connexions} \inserted{stratifications} sur $\underline O_X$, \deleted{les connexions} \inserted{celles qui sont} intégrables. Or, si $X/k$ est lisse, \deleted{\unsure{et $k$ de car. 0|et $k$ de car. nulle}} \deleted{\unsure{sur $k$ de car. nulle}} les stratifications sur $\underline O_X$ correspondent aux \inserted{connexions à courbure nulle, i.e. aux} 1-formes \uline{fermées} $\omega$ sur $X$. \margin{\uline{ici $k$ de car. 0} !} Dire que la stratification \inserted{(ou la connexion)} est triviale \deleted{\unsure{(ou encore, la conn.|(ou connex., la conn.}} revient à dire que \uline{$\omega$ est une dérivée logarithmique}. Dire qu'elle est intégrable revient donc à dire qu'il y a un entier $n>0$ tel que $n\omega$ soit une dérivée logarithmique. Comme le groupe des dérivées logarithmiques est un groupe isomorphe à un $L=\mathbb Z^{\unsure{r|n}}$ [\unsure{Rappeler|Rappel de|d’après} Rosenlicht], il est clair
% archive sheet number: 191
